\documentclass[11pt,a4paper]{article}

\usepackage[utf8]{inputenc}
\usepackage[T1]{fontenc}
\usepackage[spanish,es-tabla,es-nodecimaldot,es-noshorthands]{babel}
\usepackage[a4paper,left=2.2cm,right=2.2cm,top=2.4cm,bottom=2.4cm,headheight=22pt,headsep=16pt]{geometry}

\usepackage{graphicx}
\usepackage{xcolor}
\usepackage{colortbl}
\usepackage{booktabs}
\usepackage{tabularx}
\usepackage{longtable}
\usepackage{array}
\usepackage{enumitem}
\usepackage{microtype}
\usepackage{titlesec}
\usepackage{fancyhdr}
\usepackage{tcolorbox}
\usepackage{caption}
\usepackage{eurosym}
\usepackage[hidelinks]{hyperref}

% Paleta corporativa de AVA (identica a los entregables de las iteraciones)
\definecolor{avaprimary}{HTML}{17324D}
\definecolor{avaaccent}{HTML}{158C8C}
\definecolor{avamuted}{HTML}{536273}
\definecolor{avabg}{HTML}{F1F6F8}
\definecolor{avaborder}{HTML}{D6E0E6}
\definecolor{avasoftteal}{HTML}{E3F2F2}

\DeclareUnicodeCharacter{20AC}{\euro}

\setlength{\parindent}{0pt}
\setlength{\parskip}{5pt}
\renewcommand{\arraystretch}{1.22}

% Estilo de secciones identico a Entregable 2
\titleformat{\section}
  {\Large\bfseries\color{avaprimary}}
  {\thesection}
  {0.8em}
  {}
  [\vspace{3pt}{\color{avaaccent}\titlerule[0.4pt]}]
\titlespacing*{\section}{0pt}{12pt}{8pt}

\titleformat{\subsection}
  {\large\bfseries\color{avaaccent}}
  {\thesubsection}
  {0.7em}
  {}
\titlespacing*{\subsection}{0pt}{10pt}{4pt}

% Encabezado y pie de pagina
\pagestyle{fancy}
\fancyhf{}
\lhead{\small\color{avamuted}Plan de Gestión de Riesgos}
\rhead{\small\color{avamuted}Proyecto AUTOVIDA -- AVA}
\cfoot{\small\color{avamuted}\thepage}
\renewcommand{\headrulewidth}{0.4pt}
\renewcommand{\headrule}{\hbox to\headwidth{\color{avaborder}\leaders\hrule height \headrulewidth\hfill}}
\renewcommand{\footrulewidth}{0pt}

\captionsetup{
  font=small,
  labelfont={bf,color=avaprimary},
  textfont={color=avamuted},
  skip=5pt
}

\setlist[itemize]{leftmargin=1.3em,itemsep=2pt,topsep=2pt}

\newtcolorbox{avabox}[1][]{
  colback=avabg,
  colframe=avaborder,
  boxrule=0.5pt,
  arc=4pt,
  left=11pt,
  right=11pt,
  top=9pt,
  bottom=9pt,
  #1
}

\newcolumntype{L}[1]{>{\raggedright\arraybackslash}p{#1}}
\newcolumntype{C}[1]{>{\centering\arraybackslash}p{#1}}
\newcolumntype{Y}{>{\raggedright\arraybackslash}X}

\hypersetup{
  pdftitle={Plan de Gestión de Riesgos: Proyecto AUTOVIDA (26-175) -- AVA},
  pdfauthor={Equipo DGP1\_Jueves 8,30\_5},
  pdfsubject={Dirección y Gestión de Proyectos -- Gestión de Riesgos}
}

\begin{document}

% =========================================================================
% PORTADA (pagina 1 -- mismo formato que los entregables de las iteraciones)
% =========================================================================
\begin{titlepage}
  \thispagestyle{empty}
  \vspace*{0.05cm}
  {\color{avaaccent}\rule{\linewidth}{3pt}\par}
  \vspace{0.55cm}

  \begin{center}
    \includegraphics[width=0.42\linewidth]{img/logo_ava.jpeg}\par
    \vspace{0.2cm}
    {\large\bfseries\color{avamuted} Dirección y Gestión de Proyectos\par}
    \vspace{0.08cm}
    {\normalsize\color{avamuted} Grado en Ingeniería Informática -- Universidad de Granada\par}
  \end{center}

  \vspace{0.65cm}
  {\large\bfseries\color{avaaccent} DOCUMENTACIÓN DEL PROYECTO\par}
  \vspace{0.18cm}
  {\fontsize{24.88}{29}\selectfont\bfseries\color{avaprimary} Plan de Gestión de Riesgos\par}
  \vspace{0.22cm}
  {\Large\bfseries\color{avaaccent} Identificación, Análisis y Respuesta a Riesgos\par}
  \vspace{0.2cm}
  {\Large\color{avamuted} Proyecto AUTOVIDA (26-175) -- AVA\par}

  \vspace{0.55cm}
  \begin{avabox}
    Proyecto AUTOVIDA (26-175) -- AVA (Agenda Virtual Accesible), para el alumnado del PTVAL del Colegio de Educación Especial Fundación Purísima Concepción (Granada). Este documento identifica y prioriza los 10 riesgos principales del proyecto conforme a la metodología de gestión de riesgos de la asignatura, detallando el impacto concreto en tiempo, coste o alcance y las actividades de mitigación y evitación para los 5 riesgos prioritarios.
  \end{avabox}

  \vfill

  {\large\bfseries\color{avaprimary} Equipo (DGP1\_Jueves 8,30\_5)\par}
  \vspace{0.22cm}
  \begin{tabular}{@{}L{0.48\linewidth}L{0.48\linewidth}@{}}
    David Bacas Posadas & Adam Ben Chacha \\
    Paula Cobos González & Manuel José Fernández Poyatos \\
    Alberto David Gómez Gijón & Daniel Viloria Prieto \\
  \end{tabular}

  \vspace{0.35cm}
  \begin{tabularx}{\linewidth}{@{}l X@{}}
    {\bfseries\color{avaprimary} Contenido} & Proceso de gestión, lista priorizada de 10 riesgos y respuesta a los 5 prioritarios \\
    {\bfseries\color{avaprimary} Fecha} & 20 de noviembre de 2026 \\
  \end{tabularx}

  \vspace{0.4cm}
  {\color{avaaccent}\rule{\linewidth}{1.5pt}\par}
\end{titlepage}

\setcounter{page}{2}

% =========================================================================
% 1. PROCESO Y CRITERIOS DE GESTION DE RIESGOS
% =========================================================================
\section{Proceso y criterios de gestión de riesgos}

El propósito de este plan es asegurar que los riesgos del proyecto \textbf{AVA} se identifican, analizan, priorizan y controlan de forma continua. La lista de riesgos no es estática: se elabora al inicio del proyecto y se revisa al cierre de cada iteración (en las sesiones de seguimiento y retrospectiva en Jira), actualizando probabilidades, incorporando nuevos riesgos detectados y cerrando los superados.

\begin{itemize}
  \item \textbf{Categorías analizadas:} Riesgos de \textit{alcance y tiempo}, \textit{tecnológicos y de calidad} (Flutter, Firebase, accesibilidad real en tableta), \textit{personas y organización}, \textit{seguridad/RGPD} y \textit{dependencias externas} (ARASAAC, cuotas de Firebase).
  \item \textbf{Estrategias aplicadas:} Los 10 riesgos se ordenan de mayor a menor exposición (combinando probabilidad de ocurrencia e impacto). Los \textbf{5 riesgos prioritarios (R1--R5)} se gestionan activamente definiendo tanto medidas para \textbf{evitarlos} (reorganizar el trabajo para que no ocurran) como para \textbf{mitigarlos} (acciones inmediatas o de contingencia para reducir su efecto si aparecen). Los \textbf{5 riesgos restantes (R6--R10)}, de menor probabilidad o impacto, se \textbf{aceptan bajo monitorización} periódica.
  \item \textbf{Responsables y temporización:} El \textbf{Coordinador} de cada iteración actúa como \textit{monitor} general en Jira, con el apoyo de \textbf{Calidad} (Adam Ben Chacha) y \textbf{Usabilidad y Accesibilidad} (Paula Cobos González) en los riesgos técnicos y de producto. Las revisiones se efectúan semanalmente los jueves en clase y en el cierre de cada sprint.
\end{itemize}

% =========================================================================
% 2. LISTA PRIORIZADA DE RIESGOS DEL PROYECTO
% =========================================================================
\section{Lista de riesgos del proyecto}

La Tabla~\ref{tab:riesgos} recoge los \textbf{10 riesgos} identificados para AVA. Los cinco primeros (\textbf{R1--R5}, sombreados en cabecera de bloque) corresponden a los de mayor probabilidad e impacto y cuentan con actividades concretas de mitigación y evitación; los cinco últimos (\textbf{R6--R10}) permanecen bajo seguimiento sin gestión activa específica.

{\footnotesize
\setlength{\tabcolsep}{3.5pt}
\renewcommand{\arraystretch}{1.28}
\begin{longtable}{@{}L{2.7cm} C{1.0cm} L{3.55cm} L{4.05cm} L{4.05cm}@{}}
  \caption{Registro de los 10 riesgos del proyecto AVA (R1--R5 gestionados activamente; R6--R10 en monitorización).}
  \label{tab:riesgos} \\
  \toprule
  \textbf{Nombre del riesgo} & \textbf{Prob.} & \textbf{Impacto (tiempo / costes / alcance)} & \textbf{Actividades para mitigar el riesgo} & \textbf{Actividades para evitar el riesgo} \\
  \midrule
  \endfirsthead
  \toprule
  \textbf{Nombre del riesgo} & \textbf{Prob.} & \textbf{Impacto (tiempo / costes / alcance)} & \textbf{Actividades para mitigar el riesgo} & \textbf{Actividades para evitar el riesgo} \\
  \midrule
  \endhead
  \midrule
  \multicolumn{5}{r}{\textit{\color{avamuted}Continúa en la página siguiente}} \\
  \endfoot
  \bottomrule
  \endlastfoot

  \multicolumn{5}{@{}l}{\cellcolor{avasoftteal}\textbf{\color{avaprimary}Riesgos prioritarios gestionados activamente (R1--R5)}} \\
  \midrule

  \textbf{R1.} Sobrecarga de alcance en las iteraciones 2 y 3 (46 RF y tareas reabiertas) &
  \textbf{65 \%} &
  \textbf{Alcance y tiempo:} retraso de 1 semana en la Iteración 2 o recorte de alcance en la Iteración 3 (no se abordarían el chat accesible \texttt{RF-36--38} ni las rutinas domésticas \texttt{HU-30}). &
  Aplicar el plan de control de cambios para posponer historias de prioridad baja; dedicar 2 h extra por integrante en la última semana del sprint; redistribuir tareas entre los 6 miembros. &
  Priorizar en Jira al inicio del sprint solo historias imprescindibles; reservar un colchón de 3--4 días antes de cada entrega; revisar el tablero cada jueves en clase. \\
  \midrule

  \textbf{R2.} Fallos técnicos y curva de aprendizaje en Flutter y avisos en segundo plano (Firebase / FCM) &
  \textbf{55 \%} &
  \textbf{Tiempo y costes:} retraso de 5--7 días en el módulo de avisos (\texttt{HU-28}, \texttt{HU-29}) e incremento del 10 \% en horas de desarrollo (+28 h $\approx$ +420 €). &
  Programar en parejas las tareas críticas (\texttt{T-128--T-135}); recurrir a notificaciones locales programadas en el dispositivo como alternativa si falla Cloud Messaging. &
  Ejecutar una prueba de concepto (PoC) de notificaciones en la Semana 1 antes de integrarlas; exigir revisión por pares (\textit{pull request}) antes de fusionar en \texttt{develop}. \\
  \midrule

  \textbf{R3.} Coincidencia con entregas y exámenes de otras asignaturas &
  \textbf{50 \%} &
  \textbf{Tiempo y alcance:} caída del 25 \% en la capacidad semanal en diciembre, retrasando 4--6 días el cierre de la Iteración 3 (21/12/2026) o limitando las pruebas en tableta. &
  Adelantar el desarrollo principal a las dos primeras semanas del sprint; reasignar temporalmente tareas bloqueadas entre compañeros; realizar revisiones rápidas por Discord. &
  Cruzar el calendario de exámenes de los 6 integrantes al planificar cada sprint; repartir documentación y pruebas desde la Semana 1 para evitar acumulaciones al final. \\
  \midrule

  \textbf{R4.} Validación tardía o cambios de requisitos del Colegio Purísima Concepción &
  \textbf{40 \%} &
  \textbf{Alcance y costes:} retrabajo de 1--2 semanas en las pantallas de comedor (\texttt{HU-23--HU-25}) o permisos de dietas, incrementando el esfuerzo en un 12 \% (+35 h $\approx$ +525 €). &
  Canalizar el cambio por el control de cambios de \textit{DGP0}; incorporar solo ajustes críticos y desplazar requisitos secundarios (\texttt{RF-29--RF-32}) a una versión futura. &
  Validar los bocetos de Figma, la matriz de permisos y el esquema de datos con el profesorado (enlace con el centro) antes de comenzar la implementación de cada módulo. \\
  \midrule

  \textbf{R5.} Fallos de accesibilidad real en tabletas físicas o incidencias recurrentes de interfaz &
  \textbf{40 \%} &
  \textbf{Alcance y calidad:} rechazo de historias clave (\texttt{HU-10} voz, \texttt{HU-15} tamaño de letra, \texttt{HU-20--22} pasos y temporizador) y retraso de 4--5 días en la entrega. &
  Dedicar los últimos 3 días de cada iteración a corregir defectos visuales/táctiles; reabrir incidencias en Jira (como \texttt{T-60} y \texttt{T-81}) añadiendo pruebas de regresión. &
  Designar revisión obligatoria de Usabilidad (Paula) y Calidad (Adam) en cada PR; ejecutar \texttt{flutter test}, Accessibility Scanner, TalkBack y VoiceOver antes de cerrar cada HU. \\
  \midrule

  \multicolumn{5}{@{}l}{\cellcolor{avabg}\textbf{\color{avamuted}Riesgos secundarios bajo monitorización (R6--R10, no gestionados activamente)}} \\
  \midrule

  \textbf{R6.} Ausencia o baja temporal de un integrante por enfermedad &
  25 \% &
  \textbf{Tiempo:} pérdida de $\approx$ 12 h en un sprint, con retraso de 2--3 días en las tareas asignadas a esa persona. &
  \textit{\color{avamuted}No gestionado activamente (solo monitorización en Jira).} &
  \textit{\color{avamuted}No gestionado activamente (solo monitorización).} \\
  \midrule

  \textbf{R7.} Exposición indebida de datos sensibles de menores (perfiles y dietas) &
  20 \% &
  \textbf{Alcance y tiempo:} bloqueo del módulo de comedor (\texttt{HU-24}, \texttt{HU-25}) y 3--4 días de retrabajo en las reglas de Firestore. &
  \textit{\color{avamuted}No gestionado activamente (solo monitorización de reglas UE y datos ficticios).} &
  \textit{\color{avamuted}No gestionado activamente (solo monitorización).} \\
  \midrule

  \textbf{R8.} Caídas o límites de la API externa de pictogramas ARASAAC &
  15 \% &
  \textbf{Alcance y rendimiento:} indisponibilidad puntual en búsqueda de pictogramas (\texttt{HU-06}, \texttt{HU-15}); $\approx$ 1 día de ajuste. &
  \textit{\color{avamuted}No gestionado activamente (solo monitorización del servicio).} &
  \textit{\color{avamuted}No gestionado activamente (solo monitorización).} \\
  \midrule

  \textbf{R9.} Superación de cuotas del plan gratuito de Firebase &
  15 \% &
  \textbf{Costes:} interrupción temporal de Firestore o sobrecoste del 20--30 \% en servicios cloud (+15--30 €/mes). &
  \textit{\color{avamuted}No gestionado activamente (solo monitorización de consumo).} &
  \textit{\color{avamuted}No gestionado activamente (solo monitorización).} \\
  \midrule

  \textbf{R10.} Conflictos graves o pérdida de código en GitHub &
  15 \% &
  \textbf{Tiempo:} retraso de 1--2 días al resolver conflictos de fusión al cierre del sprint. &
  \textit{\color{avamuted}No gestionado activamente (solo monitorización de ramas).} &
  \textit{\color{avamuted}No gestionado activamente (solo monitorización).} \\
\end{longtable}
}

\newpage
% =========================================================================
% 3. SEGUIMIENTO Y REVISION
% =========================================================================
\section{Indicadores, responsables y seguimiento}

Conforme a la metodología de gestión de riesgos de la asignatura, los 5 riesgos prioritarios se monitorizan de forma semanal (sesiones de prácticas de los jueves) y al cierre de cada iteración mediante los indicadores de alerta de la Tabla~\ref{tab:indicadores}. Si algún riesgo secundario (\textbf{R6--R10}) supera el 30 \% de probabilidad o aumenta su impacto, pasará automáticamente al bloque de gestión activa en la siguiente actualización del plan.

\begin{table}[htbp]
  \centering
  \footnotesize
  \caption{Indicadores de alerta y responsables de seguimiento de los 5 riesgos prioritarios.}
  \label{tab:indicadores}
  \begin{tabularx}{\linewidth}{@{}L{1.1cm} L{3.2cm} Y L{3.4cm}@{}}
    \toprule
    \textbf{Riesgo} & \textbf{Monitor / Resp.} & \textbf{Indicador de alerta (disparador)} & \textbf{Frecuencia de revisión} \\
    \midrule
    \textbf{R1} & Coordinador & Más del 20 \% de tareas pendientes al final de la Semana 2 en Jira & Semanal (jueves) \\
    \textbf{R2} & Calidad y Equipo técnico & Bloqueo > 2 días en integración Flutter / Firebase o avisos FCM & Semanal / Pull requests \\
    \textbf{R3} & Coordinador y Moderador & Caída > 25 \% en las horas semanales registradas por integrante & Semanal (jueves) \\
    \textbf{R4} & Coordinador y Presentador & Dudas funcionales sin validar o solicitudes de cambio del centro & En cada entrega / revisión \\
    \textbf{R5} & Usabilidad (Paula) y Calidad (Adam) & Fallos en \texttt{flutter test}, TalkBack, VoiceOver o incidencias reabiertas & En cada PR y cierre de sprint \\
    \bottomrule
  \end{tabularx}
\end{table}

\end{document}
